\hypertarget{testing}{}
\hypertarget{joint-testing-and-false-discovery-control}{%
\chapter{Joint testing and false-discovery
control}\label{joint-testing-and-false-discovery-control}}

\hypertarget{the-joint-hacwald-test}{%
\section{The joint HAC/Wald test}\label{the-joint-hacwald-test}}

The null hypothesis is $H_0:\alpha=0$ for all 25 portfolios --- the joint
claim that the benchmark leaves \emph{no} portfolio with a non-zero mean
residual. The dependence-robust Wald statistic is
\begin{equation}\label{eq:ch7-wald}
W = \hat\alpha^{\top}\hat V_{\mathrm{HAC}}^{-1}\hat\alpha
\ \xrightarrow{d}\ \chi^2_{25}\quad\text{under }H_0 .
\end{equation}

\textbf{Why chi-square.} The result is the whitening argument of
\S\ref{prelim-mvn} applied to the asymptotically normal alpha vector. Under
$H_0$ and the asymptotics of \S\ref{prelim-asymp}, $\hat\alpha$ is
approximately $N(0,V_{\mathrm{HAC}})$, so the standardised vector
$z=\hat V_{\mathrm{HAC}}^{-1/2}\hat\alpha$ is approximately $N(0,I_{25})$ and
$W=z^{\top}z=\sum_{i=1}^{25}z_i^2$ is the squared norm of a standard normal
vector --- a $\chi^2$ with $25$ degrees of freedom \eqref{eq:chisq},
\eqref{eq:wald}. Crucially, nothing beyond a consistent
$\hat V_{\mathrm{HAC}}$ is required: no IID assumption, no normal residuals,
no constant covariance. That is what makes it the \emph{primary} joint
test, robust to exactly the dependence documented in Chapter~6. In the
study $W=93.9$ with $p=6.4\times10^{-10}$, so the joint zero-alpha
hypothesis is decisively rejected.

Use a linear solve rather than forming an explicit inverse:

\begin{verbatim}
W = float(alpha @ np.linalg.solve(V_hac, alpha))
p = scipy.stats.chi2.sf(W, df=len(alpha))
\end{verbatim}

\textbf{Correct conclusion.} Rejecting joint zero alpha means at least
some elements of the alpha vector are inconsistent with zero under the
model. It does not say the alphas are positive, the ranking is stable,
or a portfolio displays managerial skill.

\hypertarget{the-conventional-grs-test}{%
\section{The conventional GRS test}\label{the-conventional-grs-test}}

\begin{equation}
GRS = [(T-N-K)/N] \cdot  [1 + \bar{f}{}^{\prime}\hat{\Omega }_{f}^{-1}\bar{f}]^{-1} \cdot  \hat{\alpha }{}^{\prime}\hat{\Sigma}_{\varepsilon}^{-1}\hat{\alpha } \sim F_{N,T-N-K}
\end{equation}

Here \ensuremath{\hat{\Sigma }}\textsubscript{\ensuremath{\varepsilon }} is the N \ensuremath{\times } N contemporaneous residual
covariance, \ensuremath{\hat{\Omega }}\textsubscript{f} is the K \ensuremath{\times } K factor covariance, and \ensuremath{\bar{f}} is
the sample mean factor vector. The quadratic form
\ensuremath{\hat{\alpha }}\ensuremath{{}^{\prime}}\ensuremath{\hat{\Sigma }}\textsubscript{\ensuremath{\varepsilon }}\textsuperscript{\ensuremath{-}1}\ensuremath{\hat{\alpha }} asks whether alpha is large
relative to residual risk. The factor term discounts this evidence when
the factors themselves have a large sample mean relative to their
covariance.

\textbf{Why the $F$ distribution.} The GRS statistic is a scaled Hotelling
$T^2$. Under a correctly specified factor model with residuals that are IID
across months, jointly Gaussian, homoskedastic, and independent of the
factors, the numerator quadratic form
$\hat\alpha^{\top}\hat\Sigma_\varepsilon^{-1}\hat\alpha$ and the factor-mean
term in the denominator are \emph{independent} quadratic forms in Gaussian
vectors --- a $\chi^2_N$ scaled by a Wishart-distributed covariance. Their
ratio, scaled by $(T-N-K)/N$, is therefore exactly $F_{N,T-N-K}$ in finite
samples \eqref{eq:grs}. The denominator $1+\bar f^{\top}\hat\Omega_f^{-1}\bar f$
discounts the alpha evidence when the factors themselves carry a large
mean-variance signal (a high squared Sharpe ratio), which is the
efficiency-test interpretation: $H_0$ is equivalent to the factors spanning
the tangency portfolio. In the study $\mathrm{GRS}=3.28$ with
$p=1.2\times10^{-7}$, so it also rejects.

\textbf{Assumption ledger.} GRS is not ``wrong''; it answers a sharper
question --- exact finite-sample, not merely asymptotic --- under stronger
assumptions. But those assumptions are exactly the ones the diagnostics of
Chapter~13 reject: serial dependence, conditional heteroskedasticity, and
non-normality all break the exact-$F$ argument. GRS is therefore a
conventional sensitivity check, and the dependence-robust HAC/Wald statistic
\eqref{eq:ch7-wald} is the primary joint inference.

\hypertarget{why-25-individual-p-values-need-correction}{%
\section{Why 25 individual p-values need
correction}\label{why-25-individual-p-values-need-correction}}

If $m=25$ true nulls are each tested at level $0.05$, the expected number
of false positives is $m\times0.05=1.25$ and the probability of at least
one is $1-(1-0.05)^{25}\approx0.72$: a ``significant'' alpha found this way
is uninformative on its own. Rather than control the probability of
\emph{any} false rejection (the family-wise error rate, which Bonferroni
guards conservatively), the Benjamini--Hochberg (BH) procedure controls the
\emph{false-discovery rate} --- the expected proportion of false rejections
among all rejections,
$\mathrm{FDR}=\mathbb{E}[V/\max(R,1)]$ (\S\ref{prelim-fdr}). The procedure
is:
\begin{enumerate}
\tightlist
\item
  Sort the $p$-values $p_{(1)}\le\dots\le p_{(m)}$.
\item
  Find the largest $k$ satisfying $p_{(k)}\le kq/m$.
\item
  Reject the hypotheses with the $k$ smallest $p$-values.
\end{enumerate}
Under independence or positive dependence (PRDS) this guarantees
$\mathrm{FDR}\le(m_0/m)q\le q$, where $m_0$ is the unknown number of true
nulls \eqref{eq:bh}.

\hypertarget{a-complete-bh-example}{%
\section{A complete BH example}\label{a-complete-bh-example}}

Suppose m = 5 tests produce p-values (0.003, 0.011, 0.018, 0.040,
0.230), already sorted, and the target FDR is q = 0.05. The step-up
thresholds are:

\begin{longtable}[]{@{}p{0.14\textwidth}p{0.22\textwidth}p{0.22\textwidth}p{0.30\textwidth}@{}}
\toprule\noalign{}
Order i & p\textsubscript{(i)} & i q/m & Passes? \\
\midrule\noalign{}
\endhead
\bottomrule\noalign{}
\endlastfoot
1 & 0.003 & 0.010 & Yes \\
2 & 0.011 & 0.020 & Yes \\
3 & 0.018 & 0.030 & Yes \\
4 & 0.040 & 0.040 & Yes \\
5 & 0.230 & 0.050 & No \\
\end{longtable}

The largest passing order is $k=4$, so the first four nulls are rejected.
Notice the step-up logic: one does not stop at the first local failure
(order 5 fails, but that does not un-reject orders 1--4). The reported
adjusted $q$-values are the monotone reverse cumulative minimum
\begin{equation}\label{eq:ch7-qval}
q_{(i)} = \min_{j\ge i}\ \min\!\Bigl(1,\ \frac{m\,p_{(j)}}{j}\Bigr),
\end{equation}
which is non-decreasing in $i$ and is mapped back to the original testing
order; rejecting when $q_i\le q$ reproduces the step-up rule exactly.

BH controls FDR under independence and under broad positive-dependence
conditions. Portfolio tests are correlated, so the full dependence
structure should still be reported and the discovery family must be
stated explicitly. A more conservative arbitrary-dependence procedure
exists, but needlessly replacing BH without a design argument can
destroy power.

\begin{verbatim}
def bh_qvalues(p):
    p = np.asarray(p, float)
    m = p.size
    order = np.argsort(p)
    ranked = p[order]
    adjusted = ranked * m / np.arange(1, m + 1)
    adjusted = np.minimum.accumulate(adjusted[::-1])[::-1]
    q = np.empty(m)
    q[order] = np.clip(adjusted, 0, 1)
    return q
\end{verbatim}

\textbf{Keep separate testing families separate.} The 25 raw-alpha tests
and 25 SFA boundary tests answer different questions. Correct each
family separately and report the family definition.

In the study, five raw alphas survive 5\% BH control, but four are
significantly negative. ``Significant'' is therefore not synonymous with
``good.''

